\documentclass[10pt, letterpaper]{article}

\usepackage{amsmath}
\usepackage{geometry}
\usepackage{titlesec}
\usepackage{titling}
\usepackage{enumitem}

\title{Resume}
\author{Aiden Wenzel}

\renewcommand{\maketitle}{
	\begin{center}
		{\huge\bfseries\theauthor}
		\vspace{0.25em}

		\quad (248)-775-9622 \quad aidenwen@umich.edu \quad linkedin.com/in/aiden-wenzel \quad github.com/aiden-wenzel
	\end{center}
}


% Prevent indentation.
\setlength{\parindent}{0pt}

% Set margin rules.
\geometry{
	a4paper,
	left=0.5in,
	right=0.5in,
	bottom=0.5in,
	top=0.5in,
}

% Command for class entries.
\newcommand{\class}[1]{\textbf{#1}:}

% Custom command for section title
% Used for bar line and vertical spacing.
\newcommand{\sectionTitle}[1]{
	\section{#1}
	\hrule
	\vspace{3pt}
}

% Command for any position which has a start and end date
% \position[job position]{organization}{start date}{end date}
\newcommand{\position}[4][] {
	\subsection{#2 \small \normalfont{\textit{#1}\hfill \textit{#3 - #4}}}
}

% Section format and spacing rules.
\titleformat{\section}
{\large\bfseries}
{}
{0em}
{}

\titlespacing{\section}
{0pt} % Left margin
{1em} % Before Sep
{0pt} % After Sep

% Subsection format and spacing rules.
\titleformat{\subsection}
{\large\bfseries}
{}
{0em}
{}

\titlespacing{\subsection}
{0pt} % Left margin
{0pt} % Before Sep
{0pt} % After Sep

% Itemize rules.
\setlist[itemize]{nosep}

% Remove page numbers.
\pagenumbering{gobble}

% Start of document.
\begin{document}
\maketitle

\sectionTitle{EDUCATION}
\position[BSE, Electrical Engineering]{University of Michigan}{September 2023}{May 2027}
\begin{itemize}
	\item Computer Science Minor
	\item GPA: 3.54/4.00
\end{itemize}

\sectionTitle{SKILLS}
\begin{itemize}
	\item \textbf{Programming Languages}: C, C++, Python, Lua, MATLAB, Verilog HDL, \LaTeX. 
	\item \textbf{Linux}: Debian, WSL2, grep, cmake, make, tmux, vim, ssh.
	\item \textbf{Software}: Simulink, Waveforms Oscilloscope, LtSpice, Modelsim, Jupyter Notebooks.
\end{itemize}

\sectionTitle{PROFESSIONAL EXPERIENCE}

\position[Technical Support Specialist]{Coretek Services}{May 2024}{June 2025}
\begin{itemize}
	\item Worked with HTML, CSS, and C\# to develop intuitive ASP.NET UIs.
	\item Used C\# APIs to pull cloud resource information from MS Azure and ticket information from ServiceNow.
	\item Used KQL to query and summarize cloud resource analytics in Azure dashboards.
\end{itemize}

\position[Student]{MRI Reconstruction Lab}{Dec 2026}{Present}
\begin{itemize}
	\item Worked with a team from ECE and Radiology researching reconstruction algorithms for accelerated MRI.
	\item Reproduced reconstruction results with Python packages like scipy, PyTorch, and MIRTorch.
	\item Visually and quantitatively compared results between reconstruction methods.
	\item Presented my findings at the 2026 ECE undergraduate research symposium.
\end{itemize}

\sectionTitle{PROJECTS}
\subsection{IR Room Occupancy Counter}
\begin{itemize}
	\item Designed a robust people counting system to track the occupancy of a room.
	\item Utilized the MLX90640 24x32 IR array sensor to track the movement of heat signatures in and out of a room.
	\item Worked extensively with the MLX90640 API utilizing the I2C protocol and contour tracing algorithms.
\end{itemize}

\subsection{Fractal Visualizer}
\begin{itemize}
    \item Utilized OpenGL APIs to accurately visualize the Mandelbrot set and Julia sets.
	\item Implemented zoom functionality for near-infinite zooming into complex geometries.
\end{itemize}

\subsection{Conway's Game of Life}
\begin{itemize}
	\item Wrote custom implementation of Conway's Game of Life using SDL.
	\item Customized build system and dependency management with CMake.
\end{itemize}

\subsection{Image Compression Pipeline}
\begin{itemize}
	\item Collaborated with a team of 4 engineering students to create an image compression pipeline in Python.
	\item Utilized the Discrete Cosine Transform, Single Value Decomposition, and Huffman Encoding to compress images.
\end{itemize}

\subsection{Instrument Recognition Software}
\begin{itemize}
	\item Worked with a team of 3 other engineering students to build an instrument recognition app.
	\item Input audio data would be processed using the Librosa library. Harmonic overtones in audio samples would be extracted using the FFT algorithm.
	\item Processed audio data would be used to train a CNN which would take audio files as input, process them,
	and predict what musical instrument is featured in the sample.
\end{itemize}

\sectionTitle{RELEVANT COURSE WORK}
\begin{itemize}
	\item \class{Signals and Systems} Studied mathematical concepts such as LTI systems, convolution, Fourier series and transform, Laplace transform, sampling, and control systems. Labs included designing filters, an envelope detector, and a mock AM radio on breadboards, as well as simulating control systems in Simulink.
	\item \class{Introduction to Circuits} Studied analog circuits. Topics included nodal analysis, diodes, op amps, MOSFET transistors, RLC circuits, AC circuits, and filters.
	\item \class{Digital Logic Design} Studied combinational circuits such as multiplexers, adders, decoders, and encoders and sequential circuits such as D flip-flops, registers, multipliers, and counters. Investigated optimization and analysis techniques through timing analysis, two-level logic minimization, and state minimization. Projects involved using Modelsim and Verilog to model a trafic light controller, combinational, and sequential calculator on an FPGA.
\end{itemize}

\sectionTitle{EXTRACURRICULARS}
\position[Piccolo]{Michigan Marching Band}{August 2023}{December 2026}
\begin{itemize}
	\item Dedicated approximately 20 hours each week to rehearsals and performances in the Big House.
	\item Performed at the 2023 Big Ten Championship, 2024 ReliQuest Bowl, and 2025 Citrus Bowl. 
	\item Initiated as an official brother of Kappa Kappa Psi ($KK\Psi$), the National Honorary Band Fraternity.
    \item Performed with the Michigan Hockey Band at Yost Ice Arena.
\end{itemize}

\end{document}
